\documentclass[UTF8,11pt,a4paper]{ctexart}
\usepackage[a4paper,left=1.45cm,right=1.45cm,top=1.55cm,bottom=1.65cm]{geometry}
\usepackage{amsmath,amssymb,bm}
\usepackage{graphicx}
\usepackage{booktabs,tabularx,array,longtable}
\usepackage{enumitem}
\usepackage[dvipsnames,svgnames]{xcolor}
\usepackage{tikz}
\usetikzlibrary{arrows.meta,positioning,calc,fit,shapes.geometric}
\usepackage[most]{tcolorbox}
\usepackage{hyperref}
\usepackage{caption}
\usepackage{float}
\usepackage{fancyhdr}
\usepackage{titlesec}
\usepackage{microtype}
\IfFontExistsTF{Times New Roman}{\setmainfont{Times New Roman}}{\setmainfont{Liberation Serif}}
\IfFontExistsTF{Noto Sans CJK SC}{\setCJKmainfont{Noto Sans CJK SC}}{}
\IfFontExistsTF{Noto Sans CJK SC}{\setCJKsansfont{Noto Sans CJK SC}}{}
\IfFontExistsTF{DejaVu Sans Mono}{\setmonofont{DejaVu Sans Mono}}{\setmonofont{Latin Modern Mono}}
\definecolor{mainblue}{HTML}{0B5CAD}
\definecolor{softblue}{HTML}{EAF4FF}
\definecolor{darktext}{HTML}{1F2937}
\definecolor{softgray}{HTML}{F5F7FA}
\definecolor{accent}{HTML}{F59E0B}
\definecolor{goodgreen}{HTML}{2E7D32}
\definecolor{badred}{HTML}{B91C1C}
\definecolor{purple}{HTML}{6D5BD0}
\hypersetup{colorlinks=true,linkcolor=mainblue,urlcolor=DarkOrange,citecolor=mainblue}
\setlength{\parindent}{2em}
\setlength{\headheight}{14pt}
\setlength{\parskip}{0.22em}
\setlist[itemize]{leftmargin=2em,itemsep=0.23em,topsep=0.25em}
\setlist[enumerate]{leftmargin=2.2em,itemsep=0.23em,topsep=0.25em}
\renewcommand{\arraystretch}{1.2}
\captionsetup{font=small,labelfont=bf}
\titleformat{\section}{\Large\bfseries\sffamily\color{mainblue}}{\thesection}{0.7em}{}
\titleformat{\subsection}{\large\bfseries\sffamily\color{darktext}}{\thesubsection}{0.6em}{}
\titleformat{\subsubsection}{\normalsize\bfseries\sffamily\color{darktext}}{\thesubsubsection}{0.55em}{}
\pagestyle{fancy}
\fancyhf{}
\fancyhead[L]{\small\sffamily VARC: Vision ARC}
\fancyhead[R]{\small\sffamily \leftmark}
\fancyfoot[C]{\small\thepage}
\renewcommand{\headrulewidth}{0.4pt}
\newtcolorbox{infobox}[1]{enhanced,breakable,colback=softblue,colframe=mainblue!70!black,boxrule=0.8pt,arc=2mm,left=1.5mm,right=1.5mm,top=1mm,bottom=1mm,title=\bfseries #1}
\newtcolorbox{keybox}[1]{enhanced,breakable,colback=softgray,colframe=darktext!45,boxrule=0.6pt,arc=2mm,left=1.5mm,right=1.5mm,top=1mm,bottom=1mm,title=\bfseries #1}
\newtcolorbox{notebox}[1]{enhanced,breakable,colback=accent!8,colframe=accent!80!black,boxrule=0.7pt,arc=2mm,left=1.5mm,right=1.5mm,top=1mm,bottom=1mm,title=\bfseries #1}
\newcommand{\term}[1]{\textbf{\textcolor{mainblue}{#1}}}
\begin{document}

\section{论文基本信息}
\begin{center}
\begin{tabularx}{0.94\linewidth}{>{\bfseries}p{0.18\linewidth}X}
\toprule
题目 & VARC: Vision ARC\\
课件日期 & 2026-01-24\\
研究方向 & 视觉抽象推理 / ARC\\
一句话概括 & 把 ARC 的规则归纳从符号网格推进到视觉输入,比较基于 LLM、RNN 与视觉专用架构的抽象推理路径。\\
\bottomrule
\end{tabularx}
\end{center}

\begin{infobox}{摘要要点}
把 ARC 的规则归纳从符号网格推进到视觉输入,比较基于 LLM、RNN 与视觉专用架构的抽象推理路径。 本说明按“研究问题—核心机制—基础公式—实验阅读—局限与优化”的顺序重新组织，不保留原始材料的封面、目录和页码对应。
\end{infobox}

\section{研究问题与动机}
这项工作的出发点不是简单地替换某个网络层，而是识别现有范式中最影响\term{质量、效率、可靠性或可部署性}的瓶颈。结合论文所讨论的问题，可以将研究动机概括为以下几点：
\begin{itemize}
\item 像推理 化 言推理 图 问题转 为语 问题
\end{itemize}

\begin{notebox}{理解重点}
阅读这类论文时，应把“问题是否真实存在”“方法是否直接解决该问题”“实验是否排除了其他解释”分开判断。仅凭更大的模型、更多数据或更长训练得到的提升，不能自动视为结构创新带来的收益。
\end{notebox}

\section{核心思想与整体流程}
把 ARC 的规则归纳从符号网格推进到视觉输入,比较基于 LLM、RNN 与视觉专用架构的抽象推理路径。 从信息流角度看，其基本流程可以整理为下图。

\begin{figure}[H]
\centering
\begin{tikzpicture}[>=Latex,node distance=0.78cm]
\tikzstyle{procstep}=[draw,rounded corners,fill=softblue,align=center,minimum width=3.05cm,minimum height=0.95cm,text width=2.75cm]
\node[procstep] (a) {视觉网格};
\node[procstep,right=of a] (b) {规则归纳};
\node[procstep,right=of b] (c) {候选变换};
\node[procstep,right=of c] (d) {输出图案};
\draw[->,thick,mainblue] (a)--(b);
\draw[->,thick,mainblue] (b)--(c);
\draw[->,thick,mainblue] (c)--(d);
\end{tikzpicture}
\caption{核心信息流与处理阶段。}
\end{figure}

\subsection{方法组成与信息流}
\begin{itemize}
\item 入: 输 task token + 像 入 图 输
\item 段: 似于特定任 微 测试时训练阶 类 务 调
\item VARC (Vision ARC )核心设计
\item 离 段: 似于 线训练阶 类 预训练
\item 出:推理的 像 果 输 图 结
\item 像推理 化 言推理 图 问题转 为语 问题
\item 化:必 通 察 像 , 法 离 立存在 规则视觉 须 过观 图 总结 无 脱 视觉独
\end{itemize}

\begin{keybox}{结构解析}
\begin{itemize}
\item VARC （Vision ARC ）核心设计 ； 段 两阶 训练 ；离 段： 似于 线训练阶 类 预训练 ； 段： 似于特定任 微 测试时训练阶 类 务 调 ； 入： 输 task token + 像 入 图 输 ； 出：推理的 像 果 输 图 结
\end{itemize}
\end{keybox}

上述内容以文字方式保留方法结构，避免把整张幻灯片作为独立页面嵌入正文。

\section{数学与机制理解}
本节给出理解该工作的基础数学结构。公式用于解释方法所属范式，并不替代原论文中的完整推导。
\begin{equation}
p(y_{1:T}\mid x)=\prod_{t=1}^{T}p(y_t\mid x,y_{<t})
\end{equation}
自回归推理逐 token 生成中间步骤和答案。并行推理、思维链评估与视觉规则归纳都在重新设计这一分解方式，或检验显式文本是否忠实反映内部计算。

\begin{keybox}{从公式回到方法}
应重点观察论文究竟改变了公式中的哪一部分：是输入表示、连接矩阵、条件变量、参数化方式、训练目标，还是求解与调度过程。真正的创新通常可以在这一层被准确定位。
\end{keybox}


\section{实验结果应如何阅读}
\begin{itemize}
\item VARC (Vision ARC )核心设计
\end{itemize}

\begin{notebox}{结果解读}
\begin{itemize}
\item 实验验证
\end{itemize}
\end{notebox}

该部分以文字概括实验结论与比较条件，避免用整页截图代替分析。
实验部分不应只看单一最好数字。还需要核对模型规模、训练数据、训练步数、采样或解码预算、硬件条件以及是否使用额外蒸馏、检索或后处理。只有比较条件接近时，结果才能真正说明方法本身的优势。

\section{优点、局限与可优化空间}
\subsection{主要优点}
\begin{itemize}
\item \term{问题与方法对应明确}：论文围绕一个可识别的核心瓶颈展开，方法结构能够直接映射到该瓶颈。
\item \term{可复用的设计思想}：即使不完整复现模型，其中的分解、稀疏化、递归、条件控制或系统优化思想也可迁移到相邻任务。
\item \term{实验提供了多角度证据}：实验通常同时展示主结果、效率比较、案例或消融，使结论不完全依赖单一指标。
\end{itemize}

\subsection{局限与进一步优化}
\begin{itemize}
\item 需要进一步区分\term{结构收益}与更大算力、更多数据、不同训练技巧带来的收益，并在等参数、等计算预算下进行严格比较。
\item 对超参数、输入分布和模型规模的敏感性仍应系统分析；在一个数据集上的最优配置不一定能直接迁移。
\item 实际部署还要考虑显存峰值、并发、延迟抖动、失败案例和鲁棒性，而不仅是离线平均指标。
\item 可尝试将该方法与蒸馏、量化、缓存复用、动态路由或更强的表示学习结合，验证不同优化是否互补。
\end{itemize}

\section{结论}
总体而言，把 ARC 的规则归纳从符号网格推进到视觉输入，比较基于 LLM、RNN 与视觉专用架构的抽象推理路径。。 进一步需要注意：化:必 通 察 像 , 法 离 立存在 规则视觉 须 过观 图 总结 无 脱 视觉独;出形 固定: 出 像是否符合人 输 态 输 图 类规则。

\end{document}
